\chapter{Exemples complets — l'échelle, barreau par barreau}

Les chapitres précédents expliquaient le contrat. Celui-ci donne \textbf{du code
qui tourne}, du plus court au plus complet, pour chacune des trois choses qu'on
écrit : des règles, une IA, un plateau.

\begin{nkretenir}[Ne sautez pas de barreau]
Chaque exemple existe parce que le suivant serait incompréhensible sans lui.
\nkfile{IANegamax.cpp} fait cinq cents lignes ; il en fait cent pour qui a
d'abord lu \nkfile{IAGloutonne.cpp}, et mille pour qui ne l'a pas lu.
\end{nkretenir}

\section{Les trois échelles}

\begin{center}
\begin{tabular}{@{}p{2.3cm}p{3.7cm}p{3.7cm}p{4.1cm}@{}}
\toprule
\textbf{Écrire\ldots} & \textbf{Le plus court} & \textbf{L'intermédiaire} & \textbf{Le complet} \\
\midrule
une IA & \texttt{IAMinimale.cpp} \newline $\approx$120 l. — au hasard &
         \texttt{IAFacile.cpp} \newline $\approx$140 l. — un coup &
         \texttt{IANegamax.cpp} \newline $\approx$500 l. — négamax $\alpha\beta$ \\
\addlinespace
des règles & \texttt{RegleMinimale.cpp} \newline carré 5$\times$5 & — &
         \texttt{ConquerorRulesV2.cpp} \newline le moteur de référence \\
\addlinespace
un plateau & \texttt{mini\_3x3.json} \newline 9 cases &
         \texttt{hexagone\_6x7.json} \newline 42 cases &
         \texttt{GrilleLibre.cpp} \newline en C++, projection comprise \\
\bottomrule
\end{tabular}
\end{center}

Tous sont dans \nkfile{exemples/}, sauf le moteur de référence. Tous compilent,
et les duels ci-dessous ont été \textbf{joués}, pas supposés.

\section{D'abord : la couche confortable}

Le contrat est fait pour être \textbf{stable} — structures plates, pointeurs
bruts, aucune allocation. C'est ce qu'il faut pour une frontière binaire, et
c'est pénible à lire.

Compter « les ennemis autour de mes totems » avec le contrat nu :

\begin{nkcode}{Le contrat nu — trois boucles imbriquees}
NkcStateView v;
std::memset(&v, 0, sizeof(v));
rules->GetView(inst, st, &v);
for (uint32 i = 0; i < v.cellCount; ++i) {
    if (v.cells[i].owner != (int8)moi) continue;
    NkcCoord nb[16];
    const uint32 n = rules->GetNeighbors(inst, v.coords[i], nb, 16);
    for (uint32 k = 0; k < n; ++k)
        for (uint32 j = 0; j < v.cellCount; ++j)   // recherche a la main
            if (v.coords[j].q == nb[k].q && v.coords[j].r == nb[k].r) {
                if (v.cells[j].owner >= 0 &&
                    v.cells[j].owner != (int8)moi) ++c;
                break;
            }
}
\end{nkcode}

Avec \nkfile{Conqueror/ConquerorFacile.h} :

\begin{nkcode}{La meme chose}
Plateau p(*rules, inst, st);
for (Case c : p)
    if (c.AMoi(moi))
        contact += p.CompteVoisins(c.ou, Voisin::Ennemi, moi);
\end{nkcode}

\subsection{Ce qu'elle apporte}

\begin{center}
\begin{tabular}{@{}p{2.4cm}p{11.4cm}@{}}
\toprule
\textbf{Classe} & \textbf{À quoi ça sert} \\
\midrule
\texttt{Plateau} & lire l'état : \texttt{QuiJoue()}, \texttt{Totems(j)},
  \texttt{Avantage(moi)}, \texttt{A(coord)}, \texttt{Voisins()},
  \texttt{CompteVoisins()} — et \texttt{for (Case c : p)} \\
\texttt{Case} & une case \textbf{avec sa coordonnée} : \texttt{Vide()},
  \texttt{Bloquee()}, \texttt{AMoi(j)}, \texttt{Ennemie(j)} \\
\texttt{Coups<N>} & les coups légaux, itérables \\
\texttt{Essai} & jouer un coup \textbf{pour de faux} : clone, applique, libère
  tout seul \\
\bottomrule
\end{tabular}
\end{center}

\begin{nkretenir}[Rien de nouveau, et rien à désapprendre]
Chaque fonction appelle le contrat, et rien d'autre. Tout ce qu'on fait avec ce
fichier, on peut le faire sans — en plus long.

Tout est \textbf{inline} et \textbf{sans allocation} : aucun symbole à lier, rien
qu'un compilateur en \texttt{-O2} ne supprime. Un \texttt{Plateau} tient en
quelques pointeurs.
\end{nkretenir}

\begin{nkpiege}[Un piège quand même]
Créez l'\texttt{Essai} \textbf{hors} de la boucle des coups candidats. En créer
un par candidat coûte une allocation d'état par candidat — et c'est exactement
la boucle chaude d'une IA.
\end{nkpiege}

\subsection{La preuve : la même IA, deux fois}

\nkfile{exemples/ai/IAFacile.cpp} est \textbf{le même algorithme} que
\nkfile{IAGloutonne.cpp}, écrit avec la couche : $\approx$140 lignes contre
$\approx$240.

\begin{nkcode}{Duels joues le 2026-08-09 — 40 parties, cotes inversees}
IAFacile   vs  IAGloutonne     20 - 20     50 %   <- comportement identique
IAFacile   vs  IAMinimale      40 -  0    100 %
\end{nkcode}

Le 20--20 est le résultat \textbf{attendu} de deux joueurs identiques. C'est ce
qui prouve que la couche ne change rien au jeu — seulement à ce qu'on écrit.

\begin{nkretenir}[Par où commencer]
Lisez \nkfile{IAFacile.cpp} d'abord : sa partie qui \emph{décide} tient en quinze
lignes. Lisez \nkfile{IAGloutonne.cpp} ensuite, pour voir ce que la couche vous
épargne, et pour savoir quoi faire le jour où elle ne suffira plus.
\end{nkretenir}

\section{Une IA, en trois temps}

\subsection{Temps 1 — tirer au hasard}

Elle ne réfléchit pas. Son seul mérite est de montrer \textbf{la forme d'un
module d'IA} : neuf fonctions, dont la plupart ne font rien.

\begin{nkcode}{exemples/ai/IAMinimale.cpp — le cœur}
const uint32 total = rules->GenerateLegalMoves(inst, st, a->moves, kMaxMoves);
const uint32 n     = total < kMaxMoves ? total : kMaxMoves;
if (n == 0) return 0;          // aucun coup : l'atelier jouera PASSER
out->move = a->moves[Rand(a) % n];
return 1;
\end{nkcode}

Retenez surtout ce qu'elle \textbf{ne fait pas} : elle ne construit pas de coup.
Elle \textbf{demande} les coups légaux au moteur et en choisit un. Une IA qui
fabrique ses propres coups produit tôt ou tard un coup illégal.

\subsection{Temps 2 — regarder un coup}

Le vrai premier pas. L'algorithme tient en quatre lignes :

\begin{nkcode}{L'algorithme glouton, en entier}
pour chaque coup legal :
    cloner l'etat, jouer le coup pour de faux
    donner une note a la position obtenue
garder le coup dont la note est la meilleure
\end{nkcode}

\begin{nkcode}{exemples/ai/IAGloutonne.cpp — jouer pour de faux}
for (uint32 i = 0; i < n; ++i) {
    // JOUER POUR DE FAUX : on clone, on applique, on note. L'etat recu
    // n'est JAMAIS modifie -- le contrat l'interdit, et l'atelier compte
    // dessus pour faire tourner plusieurs parties en parallele.
    rules->CloneState(inst, a->essai, st);
    if (!rules->ApplyMove(inst, a->essai, &a->moves[i], nullptr, nullptr))
        continue;   // le moteur a refuse : il a toujours le dernier mot

    const int32 note = Noter(rules, inst, a->essai, moi, a->poidsTotems);

    // `>` et non `>=` : a note egale on garde le PREMIER. L'ordre des
    // coups est fixe par le moteur, donc ce choix est REPRODUCTIBLE.
    if (premier || note > meilleurNote) {
        meilleurNote = note; meilleur = i; premier = false;
    }
}
\end{nkcode}

L'évaluation est volontairement pauvre — un seul critère, le nombre de totems :

\begin{nkcode}{exemples/ai/IAGloutonne.cpp — Noter}
int32 miens = 0, siens = 0;
for (uint8 p = 0; p < v.playerCount; ++p) {
    if (p == moi) miens += v.totemCount[p];
    else          siens += v.totemCount[p];
}
return (miens - siens) * poids;
\end{nkcode}

\begin{nkretenir}[Ce qu'elle ne sait pas faire, et c'est le point]
Elle ne voit pas la \textbf{réponse} de l'adversaire. Un coup qui gagne trois
totems et en offre cinq au coup suivant lui paraît excellent.

C'est exactement le manque que négamax comble. Le meilleur moyen de comprendre
pourquoi négamax existe est d'avoir d'abord vu jouer celle-ci.
\end{nkretenir}

\subsection{Temps 3 — regarder loin}

Quatre pièces, dans l'ordre où elles comptent : \textbf{évaluer},
\textbf{négamax}, \textbf{alpha-bêta}, \textbf{approfondissement itératif}.

\subsubsection{Pourquoi négamax et pas minimax}

Minimax écrit deux fonctions, une qui maximise et une qui minimise. Négamax n'en
écrit qu'une, en observant que $\min(a, b) = -\max(-a, -b)$.

La valeur d'une position pour le joueur au trait est l'opposé de sa valeur pour
l'autre. On évalue \textbf{toujours du point de vue du joueur au trait}.

\begin{nkcode}{exemples/ai/IANegamax.cpp — la boucle centrale}
// L'enfant rend SA valeur, du point de vue de SON joueur au trait. On
// la ramene au point de vue de `mover` : identique si c'est le meme
// joueur qui rejoue (coup gratuit, tour multiple), opposee sinon.
if (after.current == mover) {
    score = Negamax(ai, rules, inst, ai->work[ply], depth - 1,
                    alpha, beta, ply + 1);
} else {
    score = -Negamax(ai, rules, inst, ai->work[ply], depth - 1,
                     -beta, -alpha, ply + 1);
}

if (ai->outOfTime) return 0;
if (score > best) best = score;
if (best > alpha) alpha = best;
if (alpha >= beta) break;    // COUPURE BETA
\end{nkcode}

\begin{nkpiege}[L'erreur de signe — elle s'est produite ici, en écrivant ce fichier]
La première version évaluait toujours du point de vue du joueur \textbf{racine},
tout en négligeant à chaque changement de trait. Les deux conventions se
contredisaient.

\textbf{Rien ne plantait.} Aucun test ne tombait. L'IA cherchait simplement, très
efficacement, le \textbf{pire} coup — et perdait \textbf{16 à 20 contre un joueur
aléatoire}.

Le remède n'a pas été de corriger un signe, mais de \textbf{supprimer la
possibilité de l'erreur} : le paramètre \texttt{me} a disparu de
\texttt{Negamax}. Tant qu'une fonction porte à la fois « le joueur racine » et
« le joueur au trait », il existe un endroit où l'on peut confondre les deux, et
il suffit d'un.
\end{nkpiege}

\subsubsection{Alpha-bêta}

\texttt{alpha} est le meilleur score que le joueur au trait s'est déjà garanti ;
\texttt{beta} celui que l'adversaire s'est garanti ailleurs. Dès que
\texttt{score >= beta}, on arrête : l'adversaire ne laissera jamais la partie
arriver ici.

Le résultat est \textbf{exactement} celui du négamax nu ; seul le nombre de nœuds
change. Et cette économie dépend entièrement de l'\textbf{ordre} des coups — d'où
le tri, qui est un paramètre pour qu'on puisse mesurer ce qu'il rapporte.

\subsubsection{Approfondissement itératif}

On cherche à profondeur 1, on garde le meilleur coup. Puis 2. Puis 3.

\begin{nkcode}{exemples/ai/IANegamax.cpp — Search}
// LE POINT QUI FAIT TOUT : on ne garde le resultat d'un niveau que
// s'il a ete PARCOURU EN ENTIER. Un niveau interrompu a examine les
// premiers coups seulement, et le « meilleur » qu'il propose est le
// meilleur d'un echantillon arbitraire.
if (ai->outOfTime) break;
best      = levelBestMove;
bestScore = levelBest;
reached   = depth;
\end{nkcode}

Cela paraît du gaspillage — on refait le travail à chaque fois. Ce n'en est pas :
le dernier niveau coûte plus que tous les précédents réunis. Et c'est \textbf{la
seule façon de respecter un budget de temps} : sans lui, il faudrait deviner la
profondeur atteignable, et se tromper signifie soit figer l'interface, soit ne
rien chercher.

\subsubsection{La limite assumée}

\texttt{IANegamax} détecte les parties à plus de deux joueurs et \textbf{le
dit}. L'identité $\min(a,b) = -\max(-a,-b)$ suppose un jeu à somme nulle
\textbf{à deux}. À trois, « le gain de l'un est la perte de l'autre » est faux :
il y a deux autres.

\subsection{Ce que ça donne, mesuré}

Banc de duel, 40 parties, côtés inversés une sur deux, budget 60 ms par coup,
moteur \texttt{ConquerorRulesV2} :

\begin{nkcode}{Duels joues le 2026-08-08}
Gloutonne  vs  IAMinimale (aleatoire)     40 - 0     100 %
Negamax    vs  Gloutonne                  40 - 0     100 %
Negamax    vs  IAMinimale                 40 - 0     100 %

controles :
Gloutonne  vs  Gloutonne                  20 - 20     50 %
IAMinimale vs  IAMinimale                 20 - 20     50 %
\end{nkcode}

\begin{nkretenir}[Pourquoi les deux dernières lignes comptent plus que les trois premières]
Un banc qui donne toujours 40--0 est un banc cassé. Les contrôles à 20--20
prouvent qu'il n'y a \textbf{ni biais de siège, ni bug de comptage} — et que les
100 \% au-dessus mesurent bien les IA.

C'est le premier réflexe à prendre : avant de croire un chiffre, faire tourner le
cas où l'on connaît déjà la réponse.
\end{nkretenir}

\section{Des règles : trois fonctions, ou vingt et une}

\texttt{NkcRulesVTable} compte \textbf{vingt et une} fonctions. Sur ces vingt et
une, \textbf{trois} contiennent votre jeu :

\begin{center}
\begin{tabular}{@{}p{4.0cm}p{9.8cm}@{}}
\toprule
\textbf{Fonction} & \textbf{Ce qu'elle dit} \\
\midrule
\texttt{Construire} & à quoi la partie ressemble au départ \\
\texttt{CoupsPossibles} & ce qu'on a le droit de faire \\
\texttt{Appliquer} & ce qui se passe quand on le fait \\
\bottomrule
\end{tabular}
\end{center}

Les dix-huit autres sont les mêmes pour tout le monde.
\nkfile{ConquerorRegleFacile.h} les écrit une fois : \texttt{RegleMinimale.cpp}
fait $\approx$560 lignes, \texttt{RegleFacile.cpp} en fait $\approx$110.

\begin{nkcode}{exemples/rules/RegleFacile.cpp — le jeu entier}
struct RegleFacile {
    void Construire(Grille &g) {
        g.topologie = NkcTopology::Square4;
        g.nbJoueurs = 2;
        g.AjouterRectangle(5, 5);
        g.PoserAuxCoins(2);       // symetrie 4.2
    }

    void CoupsPossibles(const Partie &p, ListeCoups &out) {
        NkcCoord voisins[8];
        for (int32 i = 0; i < p.nbCases; ++i) {
            if (p.cases[i].owner != (int8)p.joueur) continue;
            const int32 n = p.Voisins(p.ou[i], voisins, 8);
            for (int32 k = 0; k < n; ++k)
                if (p.Vide(voisins[k]))
                    out.Dupliquer(p.joueur, p.ou[i], voisins[k]);
        }
    }

    void Appliquer(Partie &p, const NkcMove &m, Evenements &ev) { /* ... */ }
};
NKC_REGLES(RegleFacile, "RegleFacile", "1.0.0", "Cours ConquerorLab")
\end{nkcode}

\subsection{Pourquoi les dix-huit autres deviennent gratuites}

Parce que \textbf{l'état cesse d'être libre}. Le contrat dit « le module choisit
sa représentation interne » — c'est sa force, et c'est ce qui coûte cher : tant
que la structure est inconnue, personne ne peut écrire sa sérialisation ni son
empreinte à votre place.

\texttt{Partie} est une structure \textbf{fixe, plate, sans pointeur}. Du coup :

\begin{center}
\begin{tabular}{@{}p{4.4cm}p{9.4cm}@{}}
\toprule
\textbf{Fonction} & \textbf{Devient} \\
\midrule
\texttt{SerializeState} & un \texttt{memcpy} \\
\texttt{CloneState} & une affectation \\
\texttt{HashState} & FNV-1a sur les octets — identique sur toute plateforme \\
\texttt{IsLegalMove} & énumérer et comparer \\
\bottomrule
\end{tabular}
\end{center}

Et ces quatre-là sont correctes \textbf{par construction}, pas par relecture.

\begin{nkretenir}[Deux pièges classiques disparaissent]
\texttt{ListeCoups::Dupliquer} fait le \texttt{memset} du coup avant de le
remplir — l'erreur numéro un des premiers modules, puisque le contrat compare les
coups \textbf{octet par octet}.

Et \texttt{PASSER} est ajouté automatiquement quand la liste reste vide, jamais
autrement : c'est exactement ce qu'exige \textsc{regles §13}.
\end{nkretenir}

\begin{nkpiege}[Ce qu'on perd, et quand il faudra partir]
Le jour où votre moteur a besoin d'un état que \texttt{Partie} ne sait pas porter
— une pile de pouvoirs, un historique, un cache d'évaluation — l'échafaudage ne
suffit plus. Vous reprendrez le contrat nu, et vous saurez exactement quelles
dix-huit fonctions écrire, \textbf{parce que vous les aurez vues à l'œuvre}.

C'est un échafaudage, pas une cage. Rien n'est caché : tout est inline et lisible
dans le header.
\end{nkpiege}

\subsection{Ce que le banc d'essai en dit}

\begin{nkcode}{RegleFacile sous tests/NkcAbiHarness.cpp}
=== 13 verifications OK, 3 echecs ===
\end{nkcode}

Les trois échecs sont ceux \textbf{codés en dur sur le moteur de référence} : le
banc cherche \texttt{portee\_duplication}, « 42 cases » et « 2 totems par
joueur ». \texttt{RegleFacile} joue un 5$\times$5 avec un totem chacun. Même
profil que \texttt{RegleMinimale} (14/16) et \texttt{GrilleLibre} (13/16).

Tout ce qui porte sur le \textbf{contrat} passe — y compris le rejeu
déterministe, le bornage des paramètres et le garde-fou \texttt{max\_tours}.

\begin{nkretenir}[Une correction que ce banc a provoquée]
À la première version, l'échafaudage n'exposait \textbf{aucun paramètre} et
n'avait \textbf{aucun garde-fou} : un jeu de règles pathologique pouvait boucler
indéfiniment, et la campagne ne rendait jamais la main. Le banc l'a signalé en
deux lignes. \texttt{max\_tours} est désormais fourni d'office, et il n'est pas
optionnel.
\end{nkretenir}

\section{Des règles, sans échafaudage : RegleMinimale.cpp}


Vingt et une fonctions dans la vtable, dont la moitié tient en trois lignes. Le
chapitre 2 les parcourt une à une ; voici seulement ce qui structure le fichier.

\begin{nkcode}{exemples/rules/RegleMinimale.cpp — le plateau, en C++}
for (int32 y = 0; y < kSide; ++y)
    for (int32 x = 0; x < kSide; ++x) {
        R->coords[y * kSide + x].q = static_cast<int16>(x);
        R->coords[y * kSide + x].r = static_cast<int16>(y);
    }
\end{nkcode}

\begin{nkpiege}[Le memset sur chaque coup]
\begin{quote}\small
\texttt{NkcMove m; std::memset(\&m, 0, sizeof(m));  // comparaison octet a octet}
\end{quote}

Le contrat compare les coups \textbf{octet par octet}. Un champ non initialisé —
même inutilisé — rend le coup différent de lui-même, et \texttt{IsLegalMove} le
refuse. C'est l'erreur numéro un des premiers modules.
\end{nkpiege}

Le moteur complet, \nkfile{modules/rules/ConquerorRulesV2.cpp}, ajoute le
chargement de plateau JSON, les paramètres, la sérialisation et l'empreinte.
Lisez-le \textbf{après} avoir écrit le vôtre, pas avant.

\section{Un plateau : les deux voies}

\subsection{En JSON, le plus petit possible}

\begin{nkcode}{exemples/boards/mini\_3x3.json}
{
  "topology": "SQUARE_4",
  "cells": [[0,0],[1,0],[2,0],
            [0,1],[1,1],[2,1],
            [0,2],[1,2],[2,2]],
  "blocked": [[1,1]],
  "starts": [
    {"player": 0, "q": 0, "r": 0, "level": 0},
    {"player": 1, "q": 2, "r": 2, "level": 0}
  ],
  "min_players": 2, "max_players": 2
}
\end{nkcode}

Neuf cases, une bloquée au centre, deux totems opposés. Une partie dure quinze
secondes et \textbf{on suit chaque coup à l'œil}.

\begin{nkretenir}[C'est le plateau sur lequel on débogue]
Chercher pourquoi un moteur se trompe sur 42 cases est un travail d'archéologue.
Sur neuf, on voit.

Réflexe à prendre : dès qu'un comportement paraît étrange, \textbf{rechargez
\texttt{mini\_3x3.json}} avant de lire une ligne de code.
\end{nkretenir}

\subsection{En C++, quand la forme est une formule}

\nkfile{exemples/rules/GrilleLibre.cpp} — trois anneaux concentriques, 19
cellules. Ni hexagone, ni carré. Le chapitre 5 le détaille ; retenez les deux
entrées qui rendent la chose possible, et celle qui rend la grille
\textbf{utilisable par une IA} :

\begin{nkcode}{Les trois entrees de geometrie (ABI 3)}
int32  (*GetCellCenter)(NkcRules self, NkcCoord c, float32 *outXY);
uint32 (*GetCellShape) (NkcRules self, NkcCoord c, float32 *outXY, uint32 cap);
uint32 (*GetNeighbors) (NkcRules self, NkcCoord c, NkcCoord *out, uint32 cap);
\end{nkcode}

Sans \texttt{GetNeighbors}, une IA devine l'adjacence à partir de
\texttt{topology} — et devine faux, en silence. \texttt{IANegamax} s'en sert pour
son terme de menace, et \textbf{s'abstient} de ce terme quand elle n'est pas
déclarée : une évaluation qui devine faux est pire qu'une évaluation qui ignore
le critère.

\section{Le chemin conseillé}

\begin{enumerate}
  \item \textbf{Copiez \nkfile{IAMinimale.cpp}} dans \nkfile{travail/ai/},
        changez son nom dans \texttt{FillFactory}, sauvegardez. Vérifiez qu'elle
        apparaît dans le panneau \emph{Joueurs}.
  \item \textbf{Lisez \nkfile{IAFacile.cpp}} — quinze lignes décident, le reste
        est de la plomberie qu'on recopie.
  \item \textbf{Écrivez votre évaluation} : partez d'\nkfile{IAFacile.cpp},
        changez \texttt{Noter}, et rien d'autre. C'est là qu'est le vrai travail
        d'une IA de jeu.
  \item \textbf{Faites-la jouer contre l'aléatoire}, 40 parties. Vous devez
        gagner largement. Sinon, c'est un signe — pas une malchance.
  \item \textbf{Puis seulement}, ouvrez \nkfile{IANegamax.cpp}.
\end{enumerate}

\begin{nkexo}[1 — Le contrôle avant la mesure]
Faites jouer \texttt{IAGloutonne} contre elle-même, 100 parties. Attendu :
environ 50 \%. Si vous obtenez autre chose, cherchez pourquoi \textbf{avant}
toute autre mesure — vous venez de trouver un biais de siège ou un bug de
comptage.
\end{nkexo}

\begin{nkexo}[2 — Casser le signe, exprès]
Dans \nkfile{IANegamax.cpp}, remplacez \texttt{-Negamax(...)} par
\texttt{Negamax(...)} dans la branche du changement de joueur. Recompilez,
faites-la jouer contre l'aléatoire. Que se passe-t-il ? Combien de temps
auriez-vous mis à trouver sans connaître la réponse ?
\end{nkexo}

\begin{nkexo}[3 — Ce que rapporte le tri]
\texttt{IANegamax} expose \texttt{tri\_des\_coups}. Coupez-le, relancez une
campagne à profondeur fixe, comparez le nombre de nœuds. De combien l'alpha-bêta
est-il moins efficace sans tri ? Le chiffre vous appartient — le facteur théorique
est connu, mais ce jeu-ci n'est pas la théorie.
\end{nkexo}

\begin{nkexo}[4 — Un troisième critère]
Ajoutez à \texttt{IAGloutonne} un terme de mobilité (nombre de coups
disponibles) et mesurez. Le gain est-il réel, ou dans le bruit ? Rappel du
chapitre 6 : 40 parties ne suffisent pas à trancher un écart de cinq points.
\end{nkexo}

\begin{nkexo}[5 — Le budget]
Passez le budget de \texttt{IANegamax} de 60 ms à 5 ms. Quelle profondeur
atteint-elle encore ? Gagne-t-elle toujours contre l'aléatoire ? Vous mesurez là
le rapport entre \textbf{temps de réflexion et force}, qui est la seule chose que
règle vraiment un palier de difficulté.
\end{nkexo}
